\section{Conclusions}\label{sec:conclusion}

We implemented and evaluated ASH-LR, a hybrid method for ill-conditioned
$\ell_2$-regularized logistic regression that combines stable loss evaluation, an
initial gradient phase, and a damped inexact Newton--CG phase with a ratio test and
gradient-regularized damping $\mu_k=\sqrt{\theta_k\norm{g_k}}$. The evaluation
followed a plan fixed before the final experiments, with eight regimes, ten seeds,
five baselines and four supplementary experiments.

ASH-LR converged on every problem in double precision, including those where
gradient descent and L-BFGS failed, and stable evaluation was essential for
convergence from distant starting points. However, the pre-registered efficiency
hypothesis was not supported: in none of the five primary ill-conditioned regimes
did ASH-LR reduce the work of the best matrix-free baseline, Newton--CG, by 20\%;
it required 18\% to 127\% more data passes in four regimes and was statistically on
a par in one badly scaled regime. Direct Newton was the fastest method throughout the dimensions
studied ($d\le500$). The ablation located the causes: the gradient phase adds work
that Newton-type methods do not need from $w_0=0$, and in badly scaled problems the
accuracy of the unpreconditioned CG solve, not the damping, limits progress.

For practitioners, the results suggest a simple rule: evaluate the loss stably, use
direct Newton when $d$ is small enough to form the Hessian, and otherwise use
Newton--CG, preferably on standardized features. Future work should
(i) add a diagonal or randomized Nystr\"om preconditioner~\cite{frangella2023randomized} to the CG solver, which targets the observed
bottleneck; (ii) test damped Newton--CG without the gradient phase, which the
ablation identified as promising on badly scaled data, in a new pre-registered study;
(iii) extend the evaluation to high-dimensional sparse data, where matrix-free
methods are the only option; and (iv) analyse the convergence of the
gradient-regularized ratio-test rule with inexact solves.

\section*{Code and data availability}
The implementation, tests, experiment scripts, raw results and the pre-registered
plan are available in the project repository (see the README for the URL). Spambase is
available from the UCI Machine Learning Repository under CC BY 4.0~\cite{spambase}.
